% GENERATED FILE. Edit canonical lessons, then run npm run generate:books.
% canonical-source-hash: fnv1a64:9bee1646161ca1f5
% canonical-lessons: HI-C171-dikhana, HI-C171-dikhna, HI-C171-chipana, HI-C171-chipna, HI-C171-rukna

\chapter{To show, To be seen, To hide (something), To hide (oneself), To stop}
\label{ch:hi-to-show-to-be-seen-to-hide-something-to-hide-oneself-to-stop}
\hlchaptermodality{171}

\begin{chapteropening}
\textbf{By the end of this chapter:} \emph{I can say to show, to be seen, to hide (something), to hide (oneself) and to stop.}

Complete the last lesson of chapter 171: I can say to show, to be seen, to hide (something), to hide (oneself) and to stop.
\end{chapteropening}

\section[\texorpdfstring{\hi{दिखाना} (dikhānā)}{दिखाना (dikhānā)}]{\hi{दिखाना} (dikhānā) — to show}

\label{lesson:HI-C171-dikhana}



\begin{quote}
Before the new one: say the Hindi for to lift, then the Hindi for to tell.
\end{quote}

\subsection*{You'll want to know: \hi{दिखाना}}
\textbf{\hi{दिखाना}} — \emph{dikhānā} — \textquotedblleft{}to show\textquotedblright{}.

\begin{grammarlens}[title={What you've built}]
One more word for this chapter.
\end{grammarlens}

\subsection*{Your turn}
\begin{itemize}
\raggedright
  \item \emph{Say it:} \emph{dikhānā}
  \item \emph{Say it:} \emph{dikhānā}, once more
  \item \emph{Say it:} say \emph{dikhānā}
  \item \emph{Recall:} say the Hindi for to lift, then the Hindi for to tell, then say \emph{dikhānā} again
\end{itemize}

\begin{tcolorbox}[breakable,colback=teal!4,colframe=teal!35!black,title={Before you move on}]
What does \hi{दिखाना} mean? (\textquotedblleft{}To show\textquotedblright{}.) Say it once more.
\end{tcolorbox}
\section[\texorpdfstring{\hi{दिखना} (dikhnā)}{दिखना (dikhnā)}]{\hi{दिखना} (dikhnā) — to be seen, to appear}

\label{lesson:HI-C171-dikhna}



\begin{quote}
Before the new one: say the Hindi for to tell, then the Hindi for to show.
\end{quote}

\subsection*{You'll want to know: \hi{दिखना}}
\textbf{\hi{दिखना}} — \emph{dikhnā} — \textquotedblleft{}to be seen, to appear\textquotedblright{}.

\begin{grammarlens}[title={What you've built}]
One more word for this chapter.
\end{grammarlens}

\subsection*{Your turn}
\begin{itemize}
\raggedright
  \item \emph{Say it:} \emph{dikhnā}
  \item \emph{Say it:} \emph{dikhnā}, once more
  \item \emph{Say it:} say \emph{dikhnā}
  \item \emph{Recall:} say the Hindi for to tell, then the Hindi for to show, then say \emph{dikhnā} again
\end{itemize}

\begin{tcolorbox}[breakable,colback=teal!4,colframe=teal!35!black,title={Before you move on}]
What does \hi{दिखना} mean? (\textquotedblleft{}To be seen, to appear\textquotedblright{}.) Say it once more.
\end{tcolorbox}
\section[\texorpdfstring{\hi{छिपाना} (chipānā)}{छिपाना (chipānā)}]{\hi{छिपाना} (chipānā) — to hide (something)}

\label{lesson:HI-C171-chipana}



\begin{quote}
Before the new one: say the Hindi for to show, then the Hindi for to be seen.
\end{quote}

\subsection*{You'll want to know: \hi{छिपाना}}
\textbf{\hi{छिपाना}} — \emph{chipānā} — \textquotedblleft{}to hide (something)\textquotedblright{}.

\begin{grammarlens}[title={What you've built}]
One more word for this chapter.
\end{grammarlens}

\subsection*{Your turn}
\begin{itemize}
\raggedright
  \item \emph{Say it:} \emph{chipānā}
  \item \emph{Say it:} \emph{chipānā}, once more
  \item \emph{Say it:} say \emph{chipānā}
  \item \emph{Recall:} say the Hindi for to show, then the Hindi for to be seen, then say \emph{chipānā} again
\end{itemize}

\begin{tcolorbox}[breakable,colback=teal!4,colframe=teal!35!black,title={Before you move on}]
What does \hi{छिपाना} mean? (\textquotedblleft{}To hide (something)\textquotedblright{}.) Say it once more.
\end{tcolorbox}
\section[\texorpdfstring{\hi{छिपना} (chipnā)}{छिपना (chipnā)}]{\hi{छिपना} (chipnā) — to hide (oneself)}

\label{lesson:HI-C171-chipna}



\begin{quote}
Before the new one: say the Hindi for to be seen, then the Hindi for to hide.
\end{quote}

\subsection*{You'll want to know: \hi{छिपना}}
\textbf{\hi{छिपना}} — \emph{chipnā} — \textquotedblleft{}to hide (oneself)\textquotedblright{}.

\begin{grammarlens}[title={What you've built}]
One more word for this chapter.
\end{grammarlens}

\subsection*{Your turn}
\begin{itemize}
\raggedright
  \item \emph{Say it:} \emph{chipnā}
  \item \emph{Say it:} \emph{chipnā}, once more
  \item \emph{Say it:} say \emph{chipnā}
  \item \emph{Recall:} say the Hindi for to be seen, then the Hindi for to hide, then say \emph{chipnā} again
\end{itemize}

\begin{tcolorbox}[breakable,colback=teal!4,colframe=teal!35!black,title={Before you move on}]
What does \hi{छिपना} mean? (\textquotedblleft{}To hide (oneself)\textquotedblright{}.) Say it once more.
\end{tcolorbox}
\section[\texorpdfstring{\hi{रुकना} (ruknā)}{रुकना (ruknā)}]{\hi{रुकना} (ruknā) — to stop, to wait}

\label{lesson:HI-C171-rukna}



\begin{quote}
Before the new one: say the Hindi for to hide, then the Hindi for to hide.
\end{quote}

\subsection*{You'll want to know: \hi{रुकना}}
\textbf{\hi{रुकना}} — \emph{ruknā} — \textquotedblleft{}to stop, to wait\textquotedblright{}.

\begin{grammarlens}[title={What you've built}]
One more word for this chapter.
\end{grammarlens}

\subsection*{Your turn}
\begin{itemize}
\raggedright
  \item \emph{Say it:} \emph{ruknā}
  \item \emph{Say it:} \emph{ruknā}, once more
  \item \emph{Say it:} say \emph{ruknā}
  \item \emph{Recall:} say the Hindi for to hide, then the Hindi for to hide, then say \emph{ruknā} again
\end{itemize}

\begin{tcolorbox}[breakable,colback=teal!4,colframe=teal!35!black,title={Before you move on}]
What does \hi{रुकना} mean? (\textquotedblleft{}To stop, to wait\textquotedblright{}.) Say it once more.
\end{tcolorbox}
